\section{Experimental Results}

To evaluate the performance of the Deliveroo agent, we conducted experiments on a benchmark suite of 10 delivery environments: 5 designed for single-agent scenarios and 5 for collaborative (multi-agent) scenarios. Each configuration was run for 5 minutes under the same computational constraints. The agent's score is defined by the total value of delivered parcels, accounting for decay over time and movement cost.

\subsection{Evaluation Protocol}

For each level, the agent was initialized in a fixed position and allowed to interact with its environment using either its default deliberative mode or the optional PDDL planner. The evaluation metric is the total score achieved at the end of the 5-minute episode.

\subsection{Score Results}

Table~\ref{tab:results} summarizes the scores obtained in each level.

\begin{table}
	\centering
	\caption{Performance per level: total score, average score per parcel, and number of delivered parcels for PDDL and non-PDDL agents}
	\label{tab:results}
	\begin{tabular}{|l|ccc|ccc|}
		\toprule
		\textbf{Level} & \multicolumn{3}{c|}{\textbf{With PDDL}} & \multicolumn{3}{c|}{\textbf{Without PDDL}} \\
		\cmidrule{2-7}
		& \textbf{Total} & \textbf{Avg} & \textbf{\# Parcels} & \textbf{Total} & \textbf{Avg} & \textbf{\# Parcels} \\
		\midrule
		25c1\_1 & 330 & 10.0 & 33 & 380 & 10.0 & 38 \\
		25c1\_2 & 222 & 6.59 & 32 & 1079 & 19.74 & 53 \\
		25c1\_3 & 1104 & 13.51 & 79 & 1303 & 14.46 & 87 \\
		25c1\_4 & 1169 & 67.18 & 18 & 2569 & 77.85 & 33 \\
		25c1\_5 & 2616 & 61.57 & 39 & 3955 & 68.82 & 57 \\
		25c1\_6 & 619 & 20.64 & 28 & 839 & 22.75 & 36 \\
		\midrule
		% optional average row
	\end{tabular}
\end{table}



\subsection{Discussion}

As shown in Table~\ref{tab:results}, the collaborative mode consistently outperforms the single-agent setup. This improvement stems from the agents' ability to share parcel and opponent beliefs, split the map for exploration, and avoid redundant pickups. Notably, the gain is most significant in larger or more spatially distributed levels, where coordination reduces travel overhead and improves delivery rate.

These results validate the effectiveness of the BDI architecture in both standalone and cooperative contexts, and they demonstrate that structured agent collaboration can yield significant efficiency gains in decentralized delivery tasks.
